\documentclass[11pt,a4paper]{article}
\usepackage[margin=1in]{geometry}
\usepackage{amsmath,amssymb}
\usepackage{booktabs}
\usepackage{hyperref}
\usepackage{graphicx}

\title{What width, data and dimension each fix in the FCN3 theory--sampler comparison}
\author{}
\date{2026-11-10}

\begin{document}
\sloppy
\maketitle

The three-layer network's theory and the exact posterior (collapsed HMC) differ in $L_3$, in the read-in action and in
the second hidden layer. Three limits can be approached separately with the runs on disk: more width at fixed theory
($N_1=N_2=N$ doubled with $\chi$ doubled), more data at fixed theory (the EK limit: $P$ doubled with $\kappa$ doubled, so
$\kappa/P$ and every parameter of the theory stay the same), and larger input dimension (the $\kappa_0=0.015$ ray, $d=150$,
300, 600 with $N=1500$, $P=3000(d/150)^2$). This note collects what each one fixes. Each fixes something different, and
none fixes the number of read-in neurons in the outer well, which each run sets early and then keeps.

Common setting: $\sigma_A^2=8/1500\cdot(150/d)$, $\sigma_B=1.2554$, $\sigma_w^2=0.5$, $\epsilon=-0.3\sqrt6$, runs started from the
prior unless stated, 1000 iterations with 300 of burn-in. ``Outer count'' is the number of read-in neurons beyond the
theory's barrier at that point; ``outer-well bottom'' is the most probable $|w_0|$ among them.

\section*{Width: doubling $N$ raises $L_3$}

On the ray at $d=600$, $P=48000$, the runs at $N=1500$ and $N=3000$ ($\chi=2$) use the same dataset.
\begin{center}\small
\begin{tabular}{@{}lccc@{}}
\toprule
& outer count (fraction) & outer-well bottom $|w_0|$ & $L_3$\\
\midrule
theory & 1.73\% & 0.898 & 0.720\\
$N=1500$, started at the theory's count & 26 (1.73\%) & 0.89 & 0.686\\
$N=1500$, prior and aligned starts & 32, 38 (2.1--2.5\%) & 0.85--0.87 & 0.692, 0.696\\
$N=3000$, prior and aligned starts & 75, 80 (2.5--2.7\%) & 0.83--0.85 & 0.705, 0.707\\
\bottomrule
\end{tabular}
\end{center}
Doubling the width raises $L_3$ by about 0.01, half the remaining gap, also between runs with the same outer fraction
(2.5\%: 0.696 against 0.705). It lowers the test error from 0.10--0.11 to 0.093--0.095 and the perpendicular part of
the network's output (its mean square after removing the best fit by functions of $x_0$) from 0.015 to 0.011. It
does not move the outer well towards the theory's position, because the $N=3000$ runs kept even more outer neurons.

In the sampler $L_3$ rises with the outer count (0.686, 0.692, 0.696 at 26, 32, 38 neurons), so comparisons of $L_3$
across widths need equal counts. At $d=150$, $P=3000$ the only width check showed no change in $L_3$, but those
$N=3000$ runs stopped at iterations 415 and 355 of 1000; there is no width test yet at strong learnability at $d=150$.

\section*{Data: the EK limit removes the overfitting features}

Theory held at an onset point ($d=150$, $N=1500$, $\kappa=0.015$) while $P$ is doubled twice with $\kappa\propto P$.
\begin{center}
\resizebox{\textwidth}{!}{%
\begin{tabular}{@{}lccc@{}}
\toprule
theory point & $P=3000$ (dataset 2) & $P=12000$ (dataset 0) & $P=24000$ (dataset 0)\\
\midrule
$L_3$ gap, $P\times1,2,4$ & 0.087, 0.030, 0.017 & 0.068, 0.038, 0.021 & 0.074, 0.039, 0.037\\
read-in neurons between the wells & 1.3--1.7, 0, 0.02 per network & 0.0009, 0.0001, 0 of all & 0.0001, 0, 0 of all\\
outer count (theory) & 26--27, 37, 39 (28) & 72, 96, 114 (78) & 47, 55, 60 (108)\\
far $h|\mathrm{He}_1$ group, per network & -- & 0, 0, 0 & 2.0, 4.6, 0\\
large-$h_3$ group ($|h_3|>0.25$), per network & -- & 0, 0, 0 & 0, 9, 11\\
\bottomrule
\end{tabular}}
\end{center}
More data at fixed theory closes the $L_3$ gap at the two lower points and halves it at the third. It empties the
barrier between the read-in wells: the neurons that sat there have their perpendicular weights pointed along the
training set's residual direction $\sum_\mu r_\mu x_{\perp,\mu}$ ($|\cos|$ 0.16--0.24 against 0.065 for other inner neurons),
a direction that is zero in the population, and with more data they settle in one well or the other. It removes the
far $h|\mathrm{He}_1$ group. It does not bring the outer count to the theory's (it overshoots at the lowest point and
stays at about half at the highest), and at the highest point a group of hidden neurons with large $h_3$ appears and
grows, which the theory does not have; holding the theory at the runs' outer count does not produce it either. At theory
point 12000 the hidden-layer widths drift slightly above the theory as $P$ grows (sd of $h_1$: 0.165, 0.176, 0.182
against 0.168).

\section*{Dimension: the outer well's position follows the count}

Along the ray the bottom of the sampler's outer read-in well lies 3\%, 6\% and 6\% inside the theory's at $d=150$, 300
and 600: larger $d$ alone does not close the gap. Runs started with the theory's outer count (27 at $d=300$, 26 at
$d=600$) put it 3\% and 1\% inside, so at the theory's count the gap does close with $d$. The offset follows the number
of outer neurons each run keeps: more outer neurons, well further in (the theory held at a lower count moves its own
well outward, 0.69 to 0.76 at $d=150$, $P=24000$). Along the ray $N$ stays at 1500 while $d$ grows, so $N/d$ falls
from 10 to 2.5 and the large-$d$ end is also the relatively narrow end.

\section*{What none of them fixes}

The outer count. Up to $P=3000$ on the onset line it settles at the theory's within 25 iterations; from $P=6000$ it
does not: at 6000 and 12000 it falls through the theory's value and is still falling at the last iteration, at 24000 it
starts at half the theory's and creeps up, at 48000 it is fixed. A read-in neuron changes well only if its incoming
weights and its outgoing column change together, so neither HMC's local moves nor Langevin dynamics (started from HMC
states at $P=3000$: $4\times10^5$ steps, no neuron reached the inner well) relax it. Which count is right cannot be
read from runs that keep their count; it needs a run started at the theory's count, or the equilibrium distribution
of the count from a free-energy calculation.

\paragraph{Sources.} \texttt{milestones/fcn3\_collapsed\_hmc}: \texttt{FINDINGS\_onset.md} and \texttt{debugging\_theory.tex};
plots \texttt{plots/onset/ek\_actions\_d150.png}, \texttt{plots/onset\_checks/theory\_fixed\_count*\_d150.png},
\texttt{plots/ray/well\_minima\_ray.png}; runs \texttt{runs/*ray\_d600\_P48000*}, \texttt{runs/*ptheory*},
\texttt{runs/*pscan\_d150*}, \texttt{runs/*onset\_d150*}, \texttt{runs/langevin/}.

\end{document}
